\section{Formal series, boundary sums, special functions, and integrals}
\label{sec:series}

\subsection{The ordinary generating function}

Let $C(x)=\sum_{n\ge0}\Cat_nx^n$ as a formal series.  Segner's recurrence is
exactly
\[
 C(x)=1+xC(x)^2.                                                  \tag{4.1}\label{eq:Cquadratic}
\]
Multiplying the proposed closed form by its conjugate shows that the unique
solution with constant term $1$ is
\[
 C(x)=\frac{1-\sqrt{1-4x}}{2x}
     =\frac{2}{1+\sqrt{1-4x}}.                                   \tag{4.2}\label{eq:Cgf}
\]
Inside \ipra, Equation~\eqref{eq:Cgf} means the polynomial relation
Equation~\eqref{eq:Cquadratic}, the initial coefficient $1$, and the recursive
coefficient definition of the square root.  No choice between two completed
infinite series is needed: coefficient uniqueness performs the choice.

The elementary bound
\(\Cat_n\le\binom{2n}{n}\le 4^n\) proves absolute convergence for
$\abs{x}<1/4$.  The boundary is also absolutely convergent, by the exact tail
identity below.  Therefore the analytic equality in Equation~\eqref{eq:Cgf}
holds for complex $x$ with $\abs{x}\le1/4$, with the square root branch having
value $1$ at zero.

\subsection{An exact telescoping boundary identity}

Put $a_n=\Cat_n/4^n$.  The ratio recurrence gives
\[
 \frac{a_{n+1}}{a_n}=\frac{2n+1}{2n+4}.
\]
Hence, with $b_n=2(n+1)a_n$,
\[
 b_n-b_{n+1}=a_n.                                                 \tag{4.3}\label{eq:telescoping-term}
\]
This yields more than the infinite sum printed on the source pages.

\begin{theorem}[Exact partial sum and primitive-recursive tail]
\label{thm:boundary-tail}
For every $N\ge0$, writing $m=N+1$,
\begin{align}
 \sum_{n=0}^{N}\frac{\Cat_n}{4^n}
   &=2-\frac{2(m+1)\Cat_m}{4^m}
    =2-\frac{2\binom{2m}{m}}{4^m},                              \tag{4.4}\label{eq:partial-boundary}\\
 0\le 2-\sum_{n=0}^{N}\frac{\Cat_n}{4^n}
   &\le \frac{2}{\sqrt{m+1}}.                                  \tag{4.5}\label{eq:tail-bound}
\end{align}
Consequently
\[
 \sum_{n=0}^{\infty}\frac{\Cat_n}{4^n}=2,
 \qquad
 \sum_{n=1}^{\infty}\frac{\Cat_n}{4^n}=1.                      \tag{4.6}\label{eq:boundary-sums}
\]
For error at most $1/q$, it is enough to take $m+1\ge4q^2$.
\end{theorem}

\begin{proof}
Equation~\eqref{eq:partial-boundary} is the finite telescope of
Equation~\eqref{eq:telescoping-term}, followed by
$(m+1)\Cat_m=\binom{2m}{m}$.

For the bound, put $u_m=\binom{2m}{m}/4^m$.  We prove
$u_m^2\le1/(m+1)$ by induction.  It is true at $m=0$, and
\[
 \frac{u_{m+1}}{u_m}=\frac{2m+1}{2m+2}.
\]
The induction step is the integer polynomial inequality
\[
 (2m+1)^2(m+2)\le4(m+1)^3,
 \qquad
 4(m+1)^3-(2m+1)^2(m+2)=3m+2.
\]
Thus the tail $2u_m$ has square at most $4/(m+1)$.  If
$m+1\ge4q^2$, nonnegativity and square comparison give $2u_m\le1/q$.
This is already a denominator-cleared \ipra{} modulus; the square-root notation
in Equation~\eqref{eq:tail-bound} is only expository.
\end{proof}

\subsection{The random-walk interpretation}

A simple symmetric walk first reaches $-1$ at time $2n+1$ exactly when its
first $2n$ steps form a Dyck path after reversing the sign convention and the
last step is downward.  There are $\Cat_n$ such step strings, each of
probability $2^{-(2n+1)}$.  Therefore the probability of first capture by time
$2N+1$ is
\[
 \sum_{n=0}^{N}\frac{\Cat_n}{2^{2n+1}}
 =1-\frac{(N+2)\Cat_{N+1}}{4^{N+1}},                              \tag{4.7}
\]
with remaining probability at most $1/\sqrt{N+2}$.  The limit is $1$ by
Theorem~\ref{thm:boundary-tail}.  Thus the recurrence of the one-dimensional
walk is not imported as an external probabilistic theorem; it is supplied
with a finite exact probability and an explicit modulus.

\subsection{The exponential generating function}

Define
\[
 F(x)=\sum_{n\ge0}\Cat_n\frac{x^n}{n!}.
\]
The factorial majorant $\Cat_n\le4^n$ gives a primitive-recursive tail bound
for every bounded complex $x$, so $F$ is entire in the constructive sense.
Multiplying Equation~\eqref{eq:ratio-cleared} by $x^n/n!$ and comparing
coefficients gives the differential equation
\[
 xF''+(2-4x)F'-2F=0,
 \qquad F(0)=F'(0)=1.                                            \tag{4.8}\label{eq:Fode}
\]
Let
\[
 I_0(2x)=\sum_{j\ge0}\frac{x^{2j}}{(j!)^2},
 \qquad
 I_1(2x)=\sum_{j\ge0}\frac{x^{2j+1}}{j!(j+1)!}.                 \tag{4.9}
\]
Termwise finite coefficient arithmetic shows that
$G(x)=e^{2x}(I_0(2x)-I_1(2x))$ satisfies Equation~\eqref{eq:Fode} and the same
initial coefficients.  The coefficient recurrence from the ODE is unique,
so
\[
 \sum_{n\ge0}\Cat_n\frac{x^n}{n!}
 =e^{2x}\bigl(I_0(2x)-I_1(2x)\bigr).                             \tag{4.10}\label{eq:egf}
\]
This is a formal coefficient proof first; convergence upgrades it to every
complex $x$.  For $\abs{x}\le A$, once $j\ge2A$ the ratio of successive terms
in either Bessel series is at most $1/4$, so four thirds of the first omitted
term is a primitive-recursive uniform tail bound.

The second Bessel identity is even simpler.  Equation~\eqref{eq:factorial}
gives
\[
 \frac{\Cat_n}{(2n)!}=\frac1{n!(n+1)!},
\]
whence, coefficient by coefficient,
\[
 \sum_{n\ge0}\Cat_n\frac{x^{2n}}{(2n)!}
 =\frac{I_1(2x)}x,                                                \tag{4.11}\label{eq:even-egf}
\]
where the value at $x=0$ is the constant coefficient $1$.

\subsection{Integral and moment representations}

\begin{theorem}[Catalan beta integrals]
For $n\ge0$,
\[
 \Cat_n=\frac1{2\pi}\int_0^4 x^n\sqrt{\frac{4-x}{x}}\,dx
 =\frac{2}{\pi}4^n\int_{-1}^{1}t^{2n}\sqrt{1-t^2}\,dt.           \tag{4.12}\label{eq:integrals}
\]
\end{theorem}

\begin{proof}
In the first integral substitute $x=4u$.  The result is
\[
 \frac{2\cdot4^n}{\pi}
 B\left(n+\frac12,\frac32\right)
 =\frac{2\cdot4^n}{\pi}
 \frac{\Gamma(n+1/2)\Gamma(3/2)}{\Gamma(n+2)}.
\]
Use $\Gamma(3/2)=\sqrt\pi/2$ and the finite half-integer gamma calculation in
Equation~\eqref{eq:gamma}.  This gives $\Cat_n$.  For the second formula, use
$x=4t^2$ on $[0,1]$ and symmetry.

For a direct constructive implementation, use the second integral.  The
integrand is continuous on the compact interval after defining its endpoint
values to be zero.  For $f_n(t)=t^{2n}\sqrt{1-t^2}$ and
$\abs{s-t}\le\delta$ on $[-1,1]$,
\[
 \abs{f_n(s)-f_n(t)}\le 2n\delta+\sqrt{2\delta}.
\]
This follows from the telescoping bound for $s^{2n}-t^{2n}$ and
$\abs{\sqrt a-\sqrt b}\le\sqrt{\abs{a-b}}$.  Hence a uniform mesh of
width at most
\[
 \delta=\frac{1}{2048q^2(n+1)^2 16^n}
\]
gives upper and lower rational Riemann sums, after rational square-root
bisection, whose contribution to the resulting Catalan interval has width
below $1/(4q)$.  Approximate $2/\pi$ by its Machin-series interval finely
enough that the product contributes below $1/(2q)$; the combined width is
then below $1/q$.  Alternatively, repeated integration by parts reduces the integral to the
Wallis recurrence and Equation~\eqref{eq:factorial}, so only rational
endpoint data remain.
\end{proof}

The first density
\[
 d\mu(x)=\frac1{2\pi}\sqrt{\frac{4-x}{x}}\,\ind{0<x<4}\,dx       \tag{4.13}\label{eq:moment-density}
\]
has mass $\Cat_0=1$ and moments $\Cat_n$.  Thus the Catalan sequence is a
compactly supported Stieltjes moment sequence on $[0,4]$.  To state the
Hausdorff version precisely, put $x=4y$ and let
\[
 d\nu(y)=\frac2\pi\sqrt{\frac{1-y}{y}}\,\ind{0<y<1}\,dy.
\]
Then $\nu$ is a probability measure on $[0,1]$ and
\[
 \int_0^1 y^n\,d\nu(y)=\frac{\Cat_n}{4^n}.
\]
Hence the normalized Catalan sequence $(\Cat_n/4^n)$ is a Hausdorff moment
sequence, which is the precise compact-interval formulation behind the source
page's ``version of the Hausdorff moment problem.''  These formulas also give
positivity of every finite Hankel quadratic form, although
Section~\ref{sec:hankel} proves the exact determinants without analysis.

The boundary sum follows from the integral by geometric summation under the
integral, but Theorem~\ref{thm:boundary-tail} is stronger for \ipra: it avoids
interchanging an integral and an infinite series and supplies a small exact
tail.
